% PRIVATE CONDITIONAL REDUCTION; NOT PROMOTED OR WIRED INTO THE PAPER.
% Actual generic cross-block invertibility is unproved.
\paragraph{An odd-isogeny character decomposition for dual Cartier forms.}
Let $\phi:A\to B$ be an odd-degree etale isogeny of abelian surfaces
over an algebraically closed field of characteristic two. Put
$n=\deg\phi$, let $M$ be a cubic principal polarization on $B$, and
let $L=\phi^*M$. Split the finite etale character algebra as
\[
\phi_*\mathcal O_A=\bigoplus_{\chi\in K^\vee}P_\chi,
\qquad V_\chi=H^0(B,M\otimes P_\chi),\qquad
H^0(A,L)=\bigoplus_\chi V_\chi.
\]
Each $V_\chi$ has dimension nine. The Cartier pairing is compatible
with the character action after the relative Frobenius twist. In
particular its trivial-character output pairs $V_\chi$ only with
$V_{\chi^{-1}}$: this condition is unaffected by whether one labels
the twisted character group through its Frobenius bijection or its
inverse. For inverse characters the actual product is a section of
$M^2$, so the cross pairing has the unambiguous type
\[
V_\chi\otimes V_{\chi^{-1}}
 \longrightarrow H^0(B^{(2)},M^{(2)}),\qquad
(s,t)\longmapsto C_{B/k}(st).
\]
Extend $\ell\in H^0(B^{(2)},M^{(2)})^*$ by zero on the other
characters of $H^0(A^{(2)},L^{(2)})$. The resulting alternating
form is block diagonal. Its trivial block is the nine-dimensional
degree-one form, and every other block is
\[
\begin{pmatrix}0&C_\chi(\ell)\\
C_\chi(\ell)^t&0\end{pmatrix}
\quad\hbox{on }V_\chi\oplus V_{\chi^{-1}}.
\]
No nontrivial character is its own inverse, since $n$ is odd.
Consequently, if the degree-one block has generic rank eight and
every nontrivial $\det C_\chi(\ell)$ is a nonzero polynomial, a
single common nonempty open of $\ell$ gives rank
$8+9(n-1)=9n-1$ on $H^0(A,L)$.

\paragraph{Exact unproved obligation.}
The required assertion is the generic invertibility of
\[
H^0(B,M\otimes P)\longrightarrow
H^0(B,M\otimes P^{-1})^*,\qquad
s\longmapsto\bigl(t\mapsto\ell(C_{B/k}(st))\bigr),
\]
for every relevant nontrivial odd-torsion $P$. The nonordinary
degree-one alternating-form theorem gives only the $P=\mathcal O$
block of rank eight; it does not prove these cross determinants.
Simplicity of $F_*M$, the nontriviality of $P$, and a generic
interpolation expectation do not by themselves establish this
claim. No such implication is asserted here. For $L=3d\Theta$
with $d$ odd, the usual maximal-isotropic quotient makes an odd
pullback of a cubic principal source available, but it does not
remove the cross-block gate. No numerical program or source
theorem was used to assert cross-block nonvanishing.
